\documentclass[12pt,a4paper]{article}
\usepackage[spanish]{babel}
\usepackage[utf8]{inputenc}
\usepackage[T1]{fontenc}
\usepackage[margin=2.5cm]{geometry}
\usepackage{setspace}
\usepackage{booktabs}
\usepackage{longtable}
\usepackage{enumitem}
\usepackage[colorlinks=true, linkcolor=blue, urlcolor=blue, citecolor=blue]{hyperref}
\usepackage[backend=biber, style=ieee, doi=true, url=true]{biblatex}
\addbibresource{SRS-references.bib}
\setstretch{1.15}
\setlength{\parindent}{0pt}
\setlength{\parskip}{6pt}
\title{Especificación de Requisitos de Software (SRS) \\ SGROAS — v0.9.0-rc}
\author{Equipo CET \\ Universidad Técnica Estatal de Quevedo — Carrera de Ingeniería de Software}
\date{24 de julio de 2026}

\begin{document}
\maketitle
\tableofcontents
\clearpage

% ============================================================
\section{Introducción}
\label{sec:introduccion}

\subsection{Propósito}
\label{sec:proposito}

El presente documento constituye la Especificación de Requisitos de Software (SRS) del
sistema \textbf{SGROAS} (Sistema de Gestión de Recursos Operativos, Administrativos y de
Seguridad), elaborada conforme a la especificación internacional \textbf{ISO/IEC/IEEE
29148:2018} \cite{iso29148}. Define los requisitos funcionales y no funcionales del sistema
en su estado candidato \texttt{v0.9.0-rc} (Tercera Entrega del Proyecto Fin de Curso, PFC),
su prioridad según el método MoSCoW, sus criterios de aceptación medibles y su método de
verificación.

\subsection{Alcance}
\label{sec:alcance}

El sistema SGROAS gestiona los recursos operativos, administrativos y de seguridad de una
institución de transporte. En el estado \texttt{v0.9.0-rc} el alcance comprende:
autenticación stateless mediante JWT en \textit{cookie} HttpOnly, CRUD de conductores y de
usuarios, caché Redis, procedimientos almacenados para reportes y agregaciones, y la
consolidación funcional heredada de la Entrega 1B. Quedan fuera del alcance de esta entrega
los módulos de facturación, pagos y la puesta en producción pública (prevista para la
Entrega Final, \texttt{v1.0.0}).

\subsection{Definiciones y abreviaturas}
\label{sec:definiciones}

\begin{description}
  \item[SGROAS] Sistema de Gestión de Recursos Operativos, Administrativos y de Seguridad.
  \item[JWT] JSON Web Token, estándar RFC 7519 de representación segura de afirmaciones.
  \item[CRUD] Create, Read, Update, Delete (crear, leer, actualizar, eliminar).
  \item[MoSCoW] Técnica de priorización: Must (obligatorio), Should (debería), Could (podría), Won't (no se hará).
  \item[ORM] Mapeo objeto-relacional (Hibernate/Spring Data JPA).
  \item[SP] \textit{Stored procedure} (procedimiento almacenado) o función SQL.
  \item[RNF] Requisito no funcional.
  \item[RF] Requisito funcional.
  \item[TTL] \textit{Time to Live}, duración máxima de una entrada de caché.
\end{description}

\subsection{Referencias}
\label{sec:referencias}

\begin{itemize}
  \item ISO/IEC/IEEE 29148:2018 — Ingeniería de requisitos \cite{iso29148}.
  \item INCOSE Guide to Writing Requirements v4 \cite{incose}.
  \item Cohn, M. \textit{User Stories Applied} (criterios INVEST) \cite{cohn2004}.
  \item Cockburn, A. \textit{Writing Effective Use Cases} \cite{cockburn2001}.
  \item RFC 7519 — JSON Web Token \cite{rfc7519}.
  \item RFC 7807 — Problem Details for HTTP APIs \cite{rfc7807}.
  \item OWASP Top 10:2021 \cite{owasp}.
  \item Robertson y Robertson. \textit{Mastering the Requirements Process} (plantilla Volere) \cite{robertson2012}.
  \item Pohl, K. \textit{Requirements Engineering} \cite{pohl2010}.
  \item Wiegers y Beatty. \textit{Software Requirements} \cite{wiegers2013}.
  \item ISO/IEC 25010:2011 — Calidad del producto software \cite{iso25010}.
\end{itemize}

\subsection{Resumen del documento}
\label{sec:resumen}

El documento se organiza en tres capítulos: la presente introducción, la descripción global
del producto (perspectiva, funciones, usuarios, restricciones, supuestos) y el capítulo de
requisitos específicos (funcionales, no funcionales y de interfaz externa). Cada requisito
indica su identificador persistente, \textit{rationale}, prioridad MoSCoW, criterio de
aceptación medible, método de verificación y trazabilidad hacia historia de usuario, caso de
uso, módulo, endpoint y evidencia empírica, en correspondencia con
\texttt{docs/trazabilidad/matriz.csv}.

% ============================================================
\section{Descripción global}
\label{sec:descripcion}

\subsection{Perspectiva del producto}
\label{sec:perspectiva}

SGROAS es una aplicación web de tres capas: frontend Angular 20, backend Spring Boot 3.2.x
(Java 21 LTS) y PostgreSQL 16, con Redis 7 como caché distribuida y orquestación mediante
Docker Compose. El backend expone una API REST documentada en OpenAPI 3.0 en \texttt{/api/docs}.
La arquitectura se detalla en los diagramas C4 (niveles 1–3) en \texttt{docs/arquitectura/} y en
las decisiones registradas en \texttt{docs/adr/}.

\subsection{Funciones del producto}
\label{sec:funciones}

\begin{itemize}
  \item Autenticación stateless con JWT en \textit{cookie} HttpOnly + Secure + SameSite=Strict.
  \item CRUD de conductores (alta, consulta, edición, eliminación lógica).
  \item CRUD de usuarios del sistema (alta, consulta, edición, eliminación lógica).
  \item Caché Redis en el listado paginado de conductores con TTL configurable externamente.
  \item Reportes y agregaciones mediante procedimientos almacenados y funciones SQL.
  \item Manejo uniforme de errores con ProblemDetails (RFC 7807).
\end{itemize}

\subsection{Características de los usuarios}
\label{sec:usuarios}

\begin{longtable}{@{}p{3cm}p{3.5cm}p{8cm}@{}}
\toprule
\textbf{Rol} & \textbf{Permiso} & \textbf{Descripción} \\
\midrule
\endhead
Administrador & Total & Gestiona usuarios, conductores y configuración del sistema. \\
Coordinador & Operativo & Gestiona conductores, rutas, vehículos y asignaciones. \\
Seguridad & Consulta & Consulta incidentes, reportes y estado de la flota. \\
\bottomrule
\end{longtable}

\subsection{Restricciones}
\label{sec:restricciones}

\begin{itemize}
  \item El esquema de base de datos se aplica exclusivamente desde \texttt{db/schema.sql},
  \texttt{db/seed.sql} y migraciones Flyway; queda prohibido \texttt{ddl-auto=update}.
  \item Toda operación que no sea CRUD elemental se implementa como procedimiento almacenado
  o función SQL (Bloque A.2.2 de la guía).
  \item Queda prohibida la concatenación de entrada de usuario en cualquier consulta SQL, JPQL
  o HQL, y el uso de SQL dinámico en procedimientos.
  \item Las imágenes de contenedor se fijan por \textit{digest} sha256.
\end{itemize}

\subsection{Supuestos y dependencias}
\label{sec:supuestos}

\begin{itemize}
  \item Se dispone de Docker Desktop para la orquestación y de Java 21 para la compilación.
  \item Las credenciales semilla (\texttt{admin@sgroas.com} / \texttt{admin123}) se documentan
  en el README y en \texttt{db/seed.sql}.
  \item La ejecución de las mediciones (k6, Lighthouse, JaCoCo) requiere el sistema levantado
  según el Makefile.
\end{itemize}

% ============================================================
\section{Requisitos específicos}
\label{sec:requisitos}

\subsection{Convenciones}
\label{sec:convenciones}

Cada requisito declara: \textbf{ID} (identificador persistente), \textbf{Tipo} (funcional /
no funcional), \textbf{Prioridad} (MoSCoW), \textbf{Rationale} (justificación y origen),
\textbf{Criterio de aceptación} (medible), \textbf{Método de verificación}
(\textit{inspection}, \textit{analysis}, \textit{demonstration}, \textit{test}) y
\textbf{Trazabilidad} (historia, caso de uso, módulo, endpoint, prueba, evidencia). Los
identificadores coinciden con la matriz \texttt{docs/trazabilidad/matriz.csv}.

\subsection{Requisitos funcionales}
\label{sec:rf}

\subsubsection{REQ-F-001 — Inicio de sesión}
\begin{itemize}
  \item \textbf{Tipo:} Funcional. \textbf{Prioridad:} Must.
  \item \textbf{Rationale:} El acceso al sistema requiere autenticación previa (origen:
  entrevistas con el personal de transporte).
  \item \textbf{Enunciado:} El sistema deberá autenticar al usuario mediante
  \texttt{POST /api/auth/login} y, tras validar credenciales, emitir un JWT firmado y
  establecer una \textit{cookie} HttpOnly + Secure + SameSite=Strict (RFC 7519).
  \item \textbf{Criterio de aceptación:} Con credenciales válidas responde \texttt{200 OK} y
  \textit{cookie} con el JWT; con credenciales inválidas responde \texttt{401 Unauthorized}
  con ProblemDetails.
  \item \textbf{Método de verificación:} \textit{test} (AuthServiceTest) y
  \textit{demonstration} (Postman).
  \item \textbf{Trazabilidad:} HU-001, CU-001, \texttt{AuthController},
  \texttt{POST /api/auth/login}, evidencia \texttt{docs/mediciones/sec/}.
\end{itemize}

\subsubsection{REQ-F-002 — Registro de usuario}
\begin{itemize}
  \item \textbf{Tipo:} Funcional. \textbf{Prioridad:} Must.
  \item \textbf{Rationale:} El administrador debe poder crear cuentas de acceso.
  \item \textbf{Enunciado:} El sistema deberá permitir registrar un usuario mediante
  \texttt{POST /api/auth/register} validando unicidad de correo.
  \item \textbf{Criterio de aceptación:} Registro exitoso responde \texttt{201 Created};
  correo duplicado responde \texttt{409 Conflict}.
  \item \textbf{Método de verificación:} \textit{test} (AuthServiceTest).
  \item \textbf{Trazabilidad:} HU-001, \texttt{AuthController}, \texttt{POST /api/auth/register}.
\end{itemize}

\subsubsection{REQ-F-003 — Renovación de token}
\begin{itemize}
  \item \textbf{Tipo:} Funcional. \textbf{Prioridad:} Must.
  \item \textbf{Rationale:} Evitar reautenticación frecuente manteniendo la seguridad.
  \item \textbf{Enunciado:} El sistema deberá renovar el JWT mediante
  \texttt{POST /api/auth/refresh} verificando el \textit{refresh token}.
  \item \textbf{Criterio de aceptación:} \textit{Refresh token} válido responde \texttt{200 OK}
  con nuevo JWT; inválido o revocado responde \texttt{401}.
  \item \textbf{Método de verificación:} \textit{test} (AuthServiceTest).
  \item \textbf{Trazabilidad:} HU-001, \texttt{AuthController}, \texttt{POST /api/auth/refresh}.
\end{itemize}

\subsubsection{REQ-F-004 — Cierre de sesión}
\begin{itemize}
  \item \textbf{Tipo:} Funcional. \textbf{Prioridad:} Must.
  \item \textbf{Rationale:} Permitir revocar la sesión activa.
  \item \textbf{Enunciado:} El sistema deberá revocar el token (JTI en lista negra Redis)
  mediante \texttt{POST /api/auth/logout}.
  \item \textbf{Criterio de aceptación:} Responde \texttt{204 No Content} y el token queda
  revocado en Redis.
  \item \textbf{Método de verificación:} \textit{test} (AuthServiceTest).
  \item \textbf{Trazabilidad:} HU-001, \texttt{AuthController}, \texttt{POST /api/auth/logout}.
\end{itemize}

\subsubsection{REQ-F-005 a REQ-F-009 — CRUD de conductores}
\begin{itemize}
  \item \textbf{Tipo:} Funcional. \textbf{Prioridad:} Must (cada uno).
  \item \textbf{Rationale:} La gestión de la flota requiere el mantenimiento de conductores
  (origen: REQ-005 de la Entrega 1A).
  \item \textbf{Enunciados:}
    \begin{description}
      \item[REQ-F-005] Listado paginado y filtrable: \texttt{GET /api/conductores}.
      \item[REQ-F-006] Consulta por id: \texttt{GET /api/conductores/\{id\}}.
      \item[REQ-F-007] Alta: \texttt{POST /api/conductores}.
      \item[REQ-F-008] Actualización: \texttt{PUT /api/conductores/\{id\}}.
      \item[REQ-F-009] Eliminación lógica: \texttt{DELETE /api/conductores/\{id\}}.
    \end{description}
  \item \textbf{Criterio de aceptación:} Cada operación responde según REST (\texttt{200},
  \texttt{201}, \texttt{204}); la cédula es única; la eliminación cambia el estado sin borrar
  la fila.
  \item \textbf{Método de verificación:} \textit{test} (ConductorServiceTest) y
  \textit{demonstration} (Postman).
  \item \textbf{Trazabilidad:} HU-002, CU-002, \texttt{ConductorController},
  evidencia \texttt{docs/mediciones/perf/} (REQ-F-005).
\end{itemize}

\subsubsection{REQ-F-010 a REQ-F-014 — CRUD de usuarios}
\begin{itemize}
  \item \textbf{Tipo:} Funcional. \textbf{Prioridad:} Must (cada uno).
  \item \textbf{Rationale:} El administrador debe gestionar las cuentas del sistema.
  \item \textbf{Enunciados:}
    \begin{description}
      \item[REQ-F-010] Listado paginado: \texttt{GET /api/usuarios}.
      \item[REQ-F-011] Consulta por id: \texttt{GET /api/usuarios/\{id\}}.
      \item[REQ-F-012] Alta: \texttt{POST /api/usuarios}.
      \item[REQ-F-013] Actualización: \texttt{PUT /api/usuarios/\{id\}}.
      \item[REQ-F-014] Eliminación lógica: \texttt{DELETE /api/usuarios/\{id\}}.
    \end{description}
  \item \textbf{Criterio de aceptación:} Respuestas REST correctas; el rol se restringe a los
  enumerados (ADMIN, COORDINADOR, SEGURIDAD).
  \item \textbf{Método de verificación:} \textit{test} (UsuarioController/Service) y
  \textit{demonstration}.
  \item \textbf{Trazabilidad:} HU-003, CU-003, \texttt{UsuarioController}.
\end{itemize}

\subsection{Requisitos no funcionales}
\label{sec:rnf}

\subsubsection{REQ-NF-001 — Cabeceras de seguridad HTTP}
\begin{itemize}
  \item \textbf{Tipo:} No funcional. \textbf{Prioridad:} Must.
  \item \textbf{Rationale:} Mitigar ataques de malconfiguración (OWASP A05).
  \item \textbf{Enunciado:} Toda respuesta HTTP incluirá \texttt{Strict-Transport-Security},
  \texttt{X-Frame-Options: DENY}, \texttt{X-Content-Type-Options: nosniff} y
  \texttt{Content-Security-Policy}.
  \item \textbf{Criterio de aceptación:} \texttt{curl -I} contra la raíz muestra las cuatro
  cabeceras.
  \item \textbf{Método de verificación:} \textit{test} y \textit{demonstration}
  (\texttt{docs/mediciones/sec/A05-*}). \textbf{Estado:} verificado.
\end{itemize}

\subsubsection{REQ-NF-002 — Cifrado en tránsito}
\begin{itemize}
  \item \textbf{Tipo:} No funcional. \textbf{Prioridad:} Must.
  \item \textbf{Rationale:} Confidencialidad e integridad (OWASP A02).
  \item \textbf{Enunciado:} La conexión utilizará TLS v1.3 con suites AEAD.
  \item \textbf{Criterio de aceptación:} \texttt{curl -v} muestra \texttt{TLSv1.3} y suite
  \texttt{TLS\_AES\_256\_GCM\_SHA384}.
  \item \textbf{Método de verificación:} \textit{demonstration}
  (\texttt{docs/mediciones/sec/A02-*}). \textbf{Estado:} verificado.
\end{itemize}

\subsubsection{REQ-NF-003 — Rendimiento del listado}
\begin{itemize}
  \item \textbf{Tipo:} No funcional. \textbf{Prioridad:} Must.
  \item \textbf{Rationale:} El listado de conductores es la consulta más frecuente (ISO/IEC
  25010, eficiencia de desempeño).
  \item \textbf{Enunciado:} El \texttt{GET /api/conductores} con caché caliente responderá
  con p95 menor a 200 ms.
  \item \textbf{Criterio de aceptación:} Tres corridas k6 (50 VUs, 30 s) con p95 medio de
  81,43 ms y 0\% de errores.
  \item \textbf{Método de verificación:} \textit{test} (k6), reporte
  \texttt{docs/mediciones/perf/REPORT.md}. \textbf{Estado:} verificado.
\end{itemize}

\subsubsection{REQ-NF-004 — Protección contra inyección}
\begin{itemize}
  \item \textbf{Tipo:} No funcional. \textbf{Prioridad:} Must.
  \item \textbf{Rationale:} Evitar inyección SQL (OWASP A03).
  \item \textbf{Enunciado:} Los payloads de entrada se parametrizan estrictamente; los
  procedimientos almacenados no construyen SQL dinámico.
  \item \textbf{Criterio de aceptación:} \texttt{POST} con payload \texttt{' OR '1'='1}
  responde \texttt{422} con ProblemDetails; el auditor
  \texttt{scripts/audit-sql-dynamic.sh} no reporta violaciones.
  \item \textbf{Método de verificación:} \textit{test} y \textit{analysis}
  (\texttt{docs/mediciones/sec/A03-*}). \textbf{Estado:} verificado.
\end{itemize}

\subsubsection{REQ-NF-005 — Limitación de intentos de login}
\begin{itemize}
  \item \textbf{Tipo:} No funcional. \textbf{Prioridad:} Must.
  \item \textbf{Rationale:} Mitigar fuerza bruta (OWASP A07).
  \item \textbf{Enunciado:} Seis intentos fallidos consecutivos desde la misma IP responderán
  \texttt{429 Too Many Requests}.
  \item \textbf{Criterio de aceptación:} La séptima petición (y siguientes) responde \texttt{429}.
  \item \textbf{Método de verificación:} \textit{demonstration}
  (\texttt{docs/mediciones/sec/A07-*}). \textbf{Estado:} verificado.
\end{itemize}

\subsubsection{REQ-NF-006 — Cobertura de pruebas}
\begin{itemize}
  \item \textbf{Tipo:} No funcional. \textbf{Prioridad:} Should (objetivo 60\,\% para la
  Tercera Entrega, 70\,\% para la Final).
  \item \textbf{Rationale:} Sostener la calidad y la tendencia creciente entre entregas.
  \item \textbf{Enunciado:} El reporte JaCoCo sobre el módulo de dominio y servicios alcanzará
  al menos 60\,\% de instrucciones en la Tercera Entrega.
  \item \textbf{Criterio de aceptación:} \texttt{jacoco.xml} con cobertura \(\ge\) 60\,\%.
  \item \textbf{Método de verificación:} \textit{test} (\texttt{docs/mediciones/jacoco/}).
  \textbf{Estado:} verificado (98,8\,\% de instrucciones y 85,4\,\% de ramas con 150 pruebas
  JUnit 5; umbral \texttt{check} de 60\,\% activo en \texttt{mvn verify}).
\end{itemize}

\subsection{Requisitos de interfaz externa}
\label{sec:interfaz}

\begin{itemize}
  \item \textbf{REQ-IF-001} (Must): La API REST se documentará en OpenAPI 3.0 accesible en
  \texttt{/api/docs}. Verificación: \textit{inspection} (Swagger UI, Figura 2 del informe).
  \item \textbf{REQ-IF-002} (Must): Los errores seguirán el formato ProblemDetails (RFC 7807)
  con \texttt{type}, \texttt{title}, \texttt{status}, \texttt{detail} e \texttt{instance}.
  Verificación: \textit{test} y \textit{demonstration} (Postman).
  \item \textbf{REQ-IF-003} (Should): El frontend se comunicará con el backend mediante
  \textit{cookie} HttpOnly (sin almacenar el JWT en \texttt{localStorage}).
\end{itemize}

\subsection{Matriz de trazabilidad}
\label{sec:matriz}

La matriz end-to-end se versiona en \texttt{docs/trazabilidad/matriz.csv} con las columnas
\texttt{id\_requisito}, \texttt{tipo}, \texttt{prioridad\_moscow},
\texttt{historia\_usuario}, \texttt{caso\_de\_uso}, \texttt{modulo\_codigo},
\texttt{endpoint\_api}, \texttt{prueba\_automatizada}, \texttt{tipo\_acceso},
\texttt{evidencia\_empirica} y \texttt{estado}. Cubre el 100\,\% de los requisitos Must
(14 funcionales + 5 no funcionales) vigentes en \texttt{v0.9.0-rc} y su validez se comprueba
en CI con \texttt{scripts/validate-traceability.sh}.

\subsection{Gestión del cambio}
\label{sec:gestion}

Las modificaciones a los requisitos se registran en \texttt{docs/requisitos/CHANGELOG-REQ.md}
siguiendo \textit{Keep a Changelog}; las decisiones de diseño que afectan requisitos Must se
documentan como ADR en \texttt{docs/adr/} (plantilla Nygard).

% ============================================================
\section*{Declaración de conformidad}
\addcontentsline{toc}{section}{Declaración de conformidad}

Cada requisito de esta especificación se redactó verificando las nueve características de
calidad para requisitos individuales y las seis para conjuntos de requisitos del INCOSE Guide
to Writing Requirements v4 \cite{incose}: Necessary, Appropriate, Unambiguous, Complete,
Singular, Feasible, Verifiable, Correct, Conforming (C1–C9) y Complete, Consistent, Feasible,
Comprehensible, Able to be validated, Correct (C10–C15).

% ============================================================
\newpage
\printbibliography[title={Referencias}]

\end{document}
